% ============================================================
% Literatur-Dateien einbinden
% ============================================================
% Empfehlung: Eine .bib-Datei pro Kapitel anlegen und hier einbinden.
% Export aus Zotero (Better BibTeX), JabRef oder Citavi.
%
% Verwendung im Text:
%   \autocite[Vgl.][S. 42]{schluessel}          → Fußnotenzitat mit Vgl.
%   \autocite[][S. 42]{schluessel}               → Fußnotenzitat ohne Vgl.
%   \autocites[Vgl.][S. 1]{q1}[Vgl.][S. 5]{q2}  → mehrere Quellen
% ============================================================

% Hauptdatei mit allen Quellen:
\addbibresource{bibliographies/resources/literatur.bib}

% Optional: kapitelweise aufteilen
% \addbibresource{bibliographies/resources/01-einleitung/01-einleitung.bib}
% \addbibresource{bibliographies/resources/02-theorie/02-01-theorie-unterkapitel.bib}
% \addbibresource{bibliographies/resources/03-methode/03-01-methode-unterkapitel.bib}
% \addbibresource{bibliographies/resources/04-ergebnisse/04-01-ergebnisse-unterkapitel.bib}
% \addbibresource{bibliographies/resources/05-diskussion/05-01-diskussion-unterkapitel}